\subsection{Haida Nation}

Surveyed the Haida Gwaii major stock assessment area from
Cumshewa Inlet south to Skincuttle Inlet.
In 2026 the Haida Fisheries Program (HFP) was contracting to complete
the herring spawn reconnaissance charter (19 days between March \nth{25} and April \nth{14})
and the herring spawn dive survey (15 days between April \nth{3} to \nth{17}).

Just prior to the start of the herring recon charter
the HFP vessel \boat{Haida Guardian} had a major mechanical issue
which resulted in the HFP contracting Leandre Vigneault
and his vessel the \boat{Ten Forward} to assist in the charter.
Leandre's experience in the recon charter assisted greatly in identifying,
tracking and (when required) conducting surface surveys.
HFP staff on the recon charter included Dan McNeill for the first 8 days and
Ben Penna for the following 11 days.
Both Dan and Ben are very experienced in herring spawn surveys on Haida Gwaii.
As a result, the Haida Gwaii herring recon charter went very well.
A total of 74,490 meters of spawn was identified and reported, and
6 of the smaller spawns were assessed using surface survey methodology.

The herring spawn dive survey began with
an abundance of spawn in upper Burnaby Strait.
Other spawns were located in inner Skincuttle Inlet, Poole Inlet and
out toward Scudder Point.
Spawns were heaviest along the East side of Huxley Island in Upper Burnaby.
By mid charter it became apparent that all spawn transects
could not be completed over the 15 charter days, so a `skip one'
pattern for survey transects was adopted on the \nth{11} day of the surveys.
Approximately 110 transects were completed over the 15 day charter.

Sounding and sampling of herring schools (prior to spawns) on Haida Gwaii
was conducted by the charter vessel \boat{Queens Reach}.
As a result, HFP staff did not have opportunities
to provide reliable observations related to fish behavior and characteristics.
During the recon and spawn charters whales were observed
as well as an abundance of sea lions and birds (similar to recent years).

Herring spawns on Haida Gwaii in 2026 appear to be similar to recent year.
Generally, most of the herring spawns rated medium to good
were concentrated in upper Burnaby Strait.
Other spawns throughout East Moresby were mostly very light to medium.
Spawns in Selwyn Inlet continue to be very minimal
compared to historical levels
(a trend which has been the case for the past 5+ years).

\begin{itemize}

\item Approximately 10,000 lbs of k'aaw was harvested
from the major stock assessment area of East Moresby
(primarily from the east side of Huxley Island Harbour in upper Burnaby Strait)
in early to mid-April.
Product was reported to be moderate to very good quality.

\item Most of the traditional k'aaw harvest was conducted by three Haida individuals:
one harvesting on behalf of the Skidegate Band,
another harvesting on behalf of the Old Massett Band, and
the third harvesting for personal use. 

\item There have been no known harvests of k'aaw
from other locations on Haida Gwaii (e.g., Skidegate Inlet).

\end{itemize}
